\documentclass[11pt,a4paper]{article}

\usepackage[margin=2.4cm]{geometry}
\usepackage[T1]{fontenc}
\usepackage[utf8]{inputenc}
\usepackage{lmodern}
\usepackage{microtype}
\usepackage{booktabs}
\usepackage{longtable}
\usepackage{array}
\usepackage{enumitem}
\usepackage{caption}
\usepackage[hidelinks]{hyperref}

\setlist{nosep}
\captionsetup{font=small,labelfont=bf}
\setlength{\parskip}{4pt}
\setlength{\parindent}{0pt}

% ==================================================================
% Numbers that depend on an engine run. Refresh this block only.
% Source: dashboard/src/data/us.json and india.json, generated 2026-08-28.
% Regenerated by reading the JSON files; do not edit by hand.
% ==================================================================
\newcommand{\wfRunDate}{2026-08-28}
\newcommand{\wfMinTrain}{252}
\newcommand{\wfStep}{21}
\newcommand{\usWFstart}{2020-10-06}
\newcommand{\usWFend}{2026-08-13}
\newcommand{\usWFperiods}{70}
\newcommand{\usWFfallbacks}{0}
\newcommand{\usWFwins}{37}
\newcommand{\usWFwinrate}{52.9}
\newcommand{\usWFport}{209.69}
\newcommand{\usWFbench}{148.63}
\newcommand{\usWFalpha}{61.06}
\newcommand{\usWFsharpe}{1.01}
\newcommand{\usBenchSharpe}{0.74}
\newcommand{\usWFvol}{15.78}
\newcommand{\usBenchVol}{16.53}
\newcommand{\usWFmaxdd}{-24.75}
\newcommand{\usBenchMaxdd}{-24.52}
\newcommand{\usWFtstat}{1.10}
\newcommand{\usWFpvalue}{0.273}
\newcommand{\usSharpeCIlo}{0.22}
\newcommand{\usSharpeCIhi}{2.01}
\newcommand{\usAlphaCIlo}{-3.05}
\newcommand{\usAlphaCIhi}{10.79}
\newcommand{\usRegimeOnMonths}{48}
\newcommand{\usRegimeOnPort}{99.97}
\newcommand{\usRegimeOnBench}{47.69}
\newcommand{\usRegimeOnAlpha}{52.28}
\newcommand{\usRegimeOnWin}{58.3}
\newcommand{\usRegimeOffMonths}{16}
\newcommand{\usRegimeOffPort}{12.64}
\newcommand{\usRegimeOffBench}{17.05}
\newcommand{\usRegimeOffAlpha}{-4.41}
\newcommand{\usRegimeOffWin}{37.5}
\newcommand{\usRegimeCrisisMonths}{6}
\newcommand{\usRegimeCrisisPort}{37.47}
\newcommand{\usRegimeCrisisBench}{43.83}
\newcommand{\usRegimeCrisisAlpha}{-6.36}
\newcommand{\usRegimeCrisisWin}{50.0}
\newcommand{\usISport}{30.52}
\newcommand{\usISbench}{20.48}
\newcommand{\usISsharpe}{2.25}
\newcommand{\usISbenchSharpe}{1.16}
\newcommand{\usDSRobserved}{1.01}
\newcommand{\usDSRbenchmark}{0.74}
\newcommand{\usDSRemax}{0.90}
\newcommand{\usDSRtrials}{37}
\newcommand{\usDSRtobs}{1470}
\newcommand{\usDSRskew}{-0.15}
\newcommand{\usDSRkurt}{3.99}
\newcommand{\usDSRpsr}{0.603}
\newcommand{\usDSRpvalue}{0.397}
\newcommand{\ffAlphaMonthly}{0.61}
\newcommand{\ffAlphaAnnual}{7.34}
\newcommand{\ffAlphaT}{2.26}
\newcommand{\ffAlphaP}{0.024}
\newcommand{\ffRsq}{0.787}
\newcommand{\ffMonths}{69}
\newcommand{\ffLags}{4}
\newcommand{\ffMktBeta}{0.823}
\newcommand{\ffMktT}{14.67}
\newcommand{\ffMktP}{0.000}
\newcommand{\ffSmbBeta}{-0.073}
\newcommand{\ffSmbT}{-0.68}
\newcommand{\ffSmbP}{0.494}
\newcommand{\ffHmlBeta}{-0.222}
\newcommand{\ffHmlT}{-1.95}
\newcommand{\ffHmlP}{0.051}
\newcommand{\ffRmwBeta}{0.331}
\newcommand{\ffRmwT}{3.40}
\newcommand{\ffRmwP}{0.001}
\newcommand{\ffCmaBeta}{0.052}
\newcommand{\ffCmaT}{0.39}
\newcommand{\ffCmaP}{0.700}
\newcommand{\ffMomBeta}{0.050}
\newcommand{\ffMomT}{0.82}
\newcommand{\ffMomP}{0.412}
\newcommand{\inWFstart}{2020-10-12}
\newcommand{\inWFend}{2026-08-14}
\newcommand{\inWFperiods}{69}
\newcommand{\inWFfallbacks}{0}
\newcommand{\inWFwins}{38}
\newcommand{\inWFwinrate}{55.1}
\newcommand{\inWFport}{142.96}
\newcommand{\inWFbench}{104.51}
\newcommand{\inWFalpha}{38.45}
\newcommand{\inWFsharpe}{0.74}
\newcommand{\inBenchSharpe}{0.49}
\newcommand{\inWFvol}{13.33}
\newcommand{\inBenchVol}{14.19}
\newcommand{\inWFmaxdd}{-15.60}
\newcommand{\inBenchMaxdd}{-17.23}
\newcommand{\inWFtstat}{1.83}
\newcommand{\inWFpvalue}{0.071}
\newcommand{\inSharpeCIlo}{-0.05}
\newcommand{\inSharpeCIhi}{1.62}
\newcommand{\inAlphaCIlo}{-0.19}
\newcommand{\inAlphaCIhi}{6.19}
\newcommand{\inRegimeOnMonths}{43}
\newcommand{\inRegimeOnPort}{49.55}
\newcommand{\inRegimeOnBench}{32.33}
\newcommand{\inRegimeOnAlpha}{17.22}
\newcommand{\inRegimeOnWin}{55.8}
\newcommand{\inRegimeOffMonths}{24}
\newcommand{\inRegimeOffPort}{59.72}
\newcommand{\inRegimeOffBench}{55.21}
\newcommand{\inRegimeOffAlpha}{4.51}
\newcommand{\inRegimeOffWin}{50.0}
\newcommand{\inRegimeCrisisMonths}{2}
\newcommand{\inRegimeCrisisPort}{1.71}
\newcommand{\inRegimeCrisisBench}{-0.44}
\newcommand{\inRegimeCrisisAlpha}{2.15}
\newcommand{\inRegimeCrisisWin}{100.0}
\newcommand{\inISport}{0.77}
\newcommand{\inISbench}{-2.17}
\newcommand{\inISsharpe}{-0.40}
\newcommand{\inISbenchSharpe}{-0.60}
\newcommand{\inDSRobserved}{0.74}
\newcommand{\inDSRbenchmark}{0.49}
\newcommand{\inDSRemax}{0.91}
\newcommand{\inDSRtrials}{37}
\newcommand{\inDSRtobs}{1449}
\newcommand{\inDSRskew}{-0.27}
\newcommand{\inDSRkurt}{2.30}
\newcommand{\inDSRpsr}{0.345}
\newcommand{\inDSRpvalue}{0.655}
\newcommand{\ledgerRows}{37}
% ==================================================================

\newcommand{\src}[1]{\footnote{\texttt{#1}}}
\newcommand{\code}[1]{\texttt{#1}}

\title{QuantShieldAI: Signal Research, Validation, and a Negative Result}
\author{Sanat Gupta}
\date{28 August 2026}

\begin{document}
\maketitle

\begin{abstract}
This report documents a systematic equity research program run between April and August 2026 on two universes: nine US large caps benchmarked to VOO and twenty NSE large caps benchmarked to the Nifty 50. Thirty-seven candidate signals, overlays and construction variants were evaluated; five signals and one weighting method are in the engine, three items are conditional and not deployed, and twenty-eight were rejected. The engine's expanding-window walk-forward (252-day minimum training window, 21-day steps, run of \wfRunDate) shows a cumulative excess return of \usWFalpha{} percentage points over VOO across \usWFperiods{} periods (\usWFstart{} to \usWFend), with a period-alpha $t$-statistic of \usWFtstat{} ($p = \usWFpvalue$) and a 10,000-draw bootstrap alpha interval of $[\usAlphaCIlo, \usAlphaCIhi]$ percentage points per year. India shows \inWFalpha{} points over \inWFperiods{} periods ($t = \inWFtstat$, $p = \inWFpvalue$). Neither is significant. The deflated Sharpe ratio on the daily walk-forward series with $N = \ledgerRows$ trials gives $p = \usDSRpvalue$ for the US and $p = \inDSRpvalue$ for India. A six-factor regression of \ffMonths{} monthly US walk-forward returns on Ken French's factors leaves an intercept of \ffAlphaMonthly\% a month ($t = \ffAlphaT$, $p = \ffAlphaP$, $R^2 = \ffRsq$), a single undeflated test on a hindsight universe. A controlled null experiment with one code path for every strategy puts the full engine at a Sharpe of 1.247 against 1.272 for naive equal weight over 107 monthly periods, with a bootstrap interval on the difference of $[-0.13, +0.19]$. Survivorship bias in the hand-picked universes is estimated at $+4.1$ and $+1.5$ percentage points per year, against annualized walk-forward alphas of about 3.9 and 3.0 points whose intervals include zero. A pre-registered intraday opening-range study on NIFTYBEES (132 sessions replayed, 43 triggers, 30.2\% win rate against a 65.1\% breakeven) lost 161.93 INR net on 2,000 INR notional, and a one-sided bootstrap of the exposure-adjusted daily delta put $P(\mathrm{mean} \le 0)$ at 1.0; its gate is closed. The system's defensible outputs are its validation machinery, its cost models, its guardrails, and a documented set of negative results.
\end{abstract}

\tableofcontents
\newpage

\section{Method}

\subsection{Universes and data}

The US universe is AAPL, GOOGL, AMZN, NVDA, JNJ, KO, BRK-B, COST and MSFT, benchmarked against VOO. The India universe is twenty NSE names across eight sectors (banks, IT exporters, consumer, industrial, telecom, pharma, finance, energy), benchmarked against the Nifty 50 index (\code{\^{}NSEI}).\src{quantshield/config.py} Macro inputs are \code{\^{}VIX}, \code{\^{}TNX}, GLD, UUP and USO for the US, and \code{\^{}INDIAVIX}, \code{\^{}NSEI}, USDINR=X and CL=F for India. The removed VIX term structure and copper/gold signals used \code{\^{}VIX3M}, HG=F and GC=F, which the engine no longer downloads. All prices are daily adjusted closes from Yahoo Finance via \code{yfinance}. The engine downloads 2,520 calendar days (about seven years) at each run, and the walk-forward starts one training year into that window. CL=F and USDINR=X print after the NSE close, so the research engine lags both by one day before any India signal or regime reads them.\src{quantshield/engine.py, \_lag\_after\_hours} The live planner is unchanged: at 09:25 IST the prior day's bars are already the latest available. Both universes were chosen by hand in 2026 from names that are large today; the consequences are quantified in Section~\ref{sec:survivorship}.

\subsection{Signal library}

Five signals are wired into the composite today.\src{quantshield/signals/composite.py}

\begin{itemize}
  \item \textbf{Momentum (12-1).} Return over the 252 trading days ending 21 days ago, rank-normalized.
  \item \textbf{Vol-adjusted momentum.} 252-day cumulative return divided by 63-day annualized volatility.
  \item \textbf{RSI(14) mean reversion.} $(50 - \mathrm{RSI}_{14})/50$, clipped to $[-1, 1]$; positive when oversold. A series with no moves over the window scores 0; an earlier version scored it as fully oversold.
  \item \textbf{SMA trend.} A score of 0.5 for price above the 200-day average, 0.3 for price above the 50-day average, 0.2 for the 50-day above the 200-day, mapped to $[-1, 1]$.
  \item \textbf{Cross-asset macro.} Market beta here is estimated over the last 126 days against the benchmark and is the \code{beta} field reported in the weight tables. US: beta-based tilts that penalize high-beta names when VIX exceeds 18, when the 10-year yield has risen more than 5\% over 63 days, or when oil has risen more than 5\% over 21 days. India: rupee depreciation favors IT exporters and penalizes consumer names, oil spikes penalize all names with an extra penalty for consumer names, and a Nifty drawdown of more than 5\% over 63 days penalizes banks.\src{quantshield/signals/cross\_asset.py}
\end{itemize}

Every signal is passed through \code{rank\_normalize} ($2\cdot\mathrm{rank}_{\mathrm{pct}} - 1$) before weighting.\src{quantshield/utils.py} Two further signals, the VIX term structure and the copper/gold ratio, were implemented and later removed from the composite as zero-contribution; their evidence is in Sections~\ref{sec:vixts} and~\ref{sec:cuau}. An earnings-surprise signal was removed earlier after walk-forward validation gave it no weight.

\subsection{Regime detector}

The regime is one of \code{risk\_on}, \code{risk\_off} or \code{crisis}, chosen by an additive score.\src{quantshield/signals/regime.py} For the US: VIX below 15 scores three points for risk-on, 15 to 20 two, 20 to 25 two for risk-off, 25 to 30 two for risk-off plus one for crisis, 30 to 35 two for crisis plus one for risk-off, above 35 three for crisis; VIX above 1.2 times its 50-day average adds one to crisis; gold up more than 5\% over a month adds two to crisis; oil up more than 10\% adds two to crisis; a rising dollar or a 30 basis point yield rise adds to risk-off. India uses India VIX thresholds of 15 and 22, rupee depreciation of 2\% and 4\%, oil above 90 USD (which blocks risk-on), and the Nifty relative to its 50-day average. The winning regime selects a signal-weight row (Table~\ref{tab:weights}).\src{quantshield/config.py} Confidence is the winning score divided by the total.

\begin{table}[htbp]
\centering\small
\caption{Composite weights by regime, normalized to one.}
\label{tab:weights}
\begin{tabular}{llrrrrr}
\toprule
Market & Regime & Momentum & Vol-adj mom. & Mean rev. & Trend & Cross-asset \\
\midrule
US & risk\_on & 0.35 & 0.25 & 0.05 & 0.20 & 0.10 \\
US & risk\_off & 0.20 & 0.20 & 0.15 & 0.20 & 0.15 \\
US & crisis & 0.10 & 0.10 & 0.25 & 0.15 & 0.20 \\
India & risk\_on & 0.30 & 0.25 & 0.05 & 0.20 & 0.20 \\
India & risk\_off & 0.20 & 0.20 & 0.15 & 0.20 & 0.25 \\
India & crisis & 0.15 & 0.10 & 0.30 & 0.15 & 0.30 \\
\bottomrule
\end{tabular}
\end{table}

\subsection{Portfolio construction}

The pipeline, as implemented in \code{quantshield/research/portfolio.py}, is:\src{quantshield/research/portfolio.py}

\begin{enumerate}
  \item Composite score $= \sum_k w_k s_k$ over the five rank-normalized signals with the regime's weights.
  \item Signal tilt $=$ percentile rank of the composite, normalized to sum to one.
  \item Base weights from Hierarchical Risk Parity (Ward linkage on $\sqrt{(1-\rho)/2}$, recursive bisection with inverse-variance cluster allocation). The original HRP paper uses single linkage; the two were not compared.\src{quantshield/risk/hrp.py}
  \item Blend: $0.5\,\mathrm{HRP} + 0.5\,\mathrm{tilt}$.
  \item Limits, iterated up to ten passes until all hold: clip to $[2\%, 40\%]$ (US) or $[1\%, 15\%]$ (India), renormalize, sector cap of 40\% (US) or 30\% (India), single-name cap of 25\% (US) or 15\% (India), portfolio beta cap of 1.5 against the benchmark (a bisection toward inverse-beta-squared weights). The cap uses its own beta: a 252-day estimate clipped to $[0, 5]$, distinct from the 126-day beta of the cross-asset signal.\src{quantshield/risk/position\_limits.py}
  \item Monthly CVaR cap of 3\% at 95\% confidence: while the empirical mean of the worst 5\% of monthly portfolio returns is below $-3\%$, move 10\% of the worst contributor's weight to the best contributor, up to twenty times. An earlier version compared the absolute value, so a tail made up entirely of gains counted as a breach; the sign is now checked.\src{quantshield/risk/cvar.py}
  \item Limits again.
\end{enumerate}

A volatility-targeting step (15\% US, 20\% India, scalar clipped to $[0.5, 1.5]$) used to sit between the blend and the limits. It was removed after the null experiment (Section~\ref{sec:null}) showed it to be exposure reduction with no risk-adjusted benefit; renormalizing to full investment afterwards also made it a no-op in the engine's own pipeline.

\subsection{Costs}

The engine walk-forward charges the US book a flat 10 basis points of one-way turnover at each rebalance. The India walk-forward charges the delivery (CNC) cost model used by the live book: securities transaction tax of 0.1\% on equity buys and sells, exchange transaction charge of 0.00297\%, SEBI charge of 0.0001\%, 18\% GST on those two, 0.015\% stamp duty on buys, and a depository charge of 15.93 INR per sell.\src{quantshield/costs.py} ETF legs pay no STT on buys and 0.001\% on sells. The intraday (MIS) schedule used by the opening-range study is brokerage of 0.03\% capped at 20 INR per order, 0.025\% STT on sells, 0.003\% stamp duty on buys, the same exchange and SEBI charges, and GST on brokerage and those two charges. The ETF intraday sell STT was 0 until \wfRunDate; Zerodha charges 0.025\% on equity ETFs such as NIFTYBEES, and the correction is recorded as a dated amendment in the gate file (Section~\ref{sec:orb}). The research scripts use a 5 basis point spread plus a square-root market impact model. The cost model was checked against a broker contract note of 2026-07-20: three delivery buys carried total charges of 3.18 INR against the model's estimate within 0.75 INR.\src{tests/test\_planner.py}

\subsection{Walk-forward}

The walk-forward is an expanding window: the first rebalance uses the first 252 trading days, every later rebalance uses all data to that date, and each test period is the next 21 trading days.\src{quantshield/research/backtest.py} Regime detection, signals, HRP and all limits are recomputed from training data only at each step. The target weights are bought on the first day of the test window and held without rebalancing for the 21 days, so they drift with prices; the portfolio's daily return is the return of that drifting book. At the next step the cost function is charged on the turnover from the drifted end-of-window weights to the new targets, and the first step charges the full purchase from cash. An earlier version held the target weights fixed on every day of the window, charged turnover between successive targets, and charged nothing for the first purchase. If any step raises, that step falls back to equal weight and is counted. Period alpha is the portfolio return minus the benchmark return over the 21-day test window, in percentage points; it is not beta-adjusted. Earlier documents described this as a rolling twelve-month window. That was wrong; the window has always been expanding with a 252-day minimum.

\subsection{Statistical tests}

\begin{itemize}
  \item \textbf{Period-alpha $t$-test.} One-sample $t$ on the vector of period alphas against zero.
  \item \textbf{Bootstrap.} 10,000 resamples of the paired 21-day portfolio and benchmark returns (seed 42). The Sharpe interval subtracts the same risk-free rate as the headline Sharpe, converted to a 21-day rate, and annualizes by $\sqrt{12}$; the alpha interval is the resampled mean excess return times twelve, in percentage points per year. An earlier version left the risk-free rate out of the interval, so it bracketed a different quantity from the point estimate.
  \item \textbf{Sharpe ratios.} Daily returns annualized by $\sqrt{252}$. The US risk-free rate is the end-of-sample 10-year yield (\code{\^{}TNX}) applied to every day in the series; India uses a fixed 6.5\% a year. The same rate is subtracted from the portfolio and the benchmark.
  \item \textbf{Deflated Sharpe ratio.} Bailey and L\'opez de Prado (2014), computed on the daily walk-forward return series. The annual Sharpe is converted to per-day units; its standard error uses the skewness and kurtosis of the series over $T$ daily observations; the expected maximum of $N$ null Sharpes with that standard error is computed, and the threshold $\mathrm{SR}^{*}$ is the larger of that expected maximum and the benchmark's Sharpe; the probability that the observed Sharpe exceeds $\mathrm{SR}^{*}$ is the probabilistic Sharpe ratio and $p$ is one minus it. The variance fed to the expected-maximum formula is the estimation variance of the observed Sharpe, which is a lower bound on the variance across trials, so the expected maximum is understated and the reported $p$ is a lower bound. $N = \ledgerRows$ is the row count of the signal ledger (\code{docs/signal\_ledger.md}). It counts candidates, not the parameter variants tried inside each study, so the deflation is a floor. An earlier output computed this on the in-sample 252-day Sharpe with $N = 17$ and mixed annual and per-period units; that figure is withdrawn.\src{quantshield/risk/deflated\_sharpe.py}
  \item \textbf{Fama-French decomposition.} A six-factor regression (market, SMB, HML, RMW, CMA, momentum) of calendar-month walk-forward excess returns on Ken French's monthly library, with Newey-West standard errors at $\lfloor n^{1/3} \rfloor$ lags; months with fewer than 15 sessions are dropped.\src{quantshield/signals/fama\_french.py} It runs for the US only. No factor library is wired for India, so no factor-adjusted figure exists there. An earlier output ran on synthetic factor series and was never reported. Results are in Section~\ref{sec:ff}.
  \item \textbf{Survivorship discount.} Any reported alpha should have 4.1 (US) or 1.5 (India) percentage points per year subtracted before interpretation (Section~\ref{sec:survivorship}).
\end{itemize}

\newpage
\section{Signal Studies}

Each study used a 70/30 in-sample/out-of-sample split on daily data unless stated, Spearman and Pearson information coefficients against 5-day and 21-day forward returns of the benchmark and of the equal-weight portfolio, rolling 63-day IC and its ratio to its standard deviation (ICIR), quintile sorts, VIX-regime splits, a turnover-based cost estimate, correlation with the existing signals, and a Bonferroni threshold over the signals tested to that point. Twenty-one-day forward returns overlap, so the effective sample is roughly one twentieth of the row count; the studies note this where it changes a conclusion.

\subsection{VIX term structure}
\label{sec:vixts}

\textbf{Hypothesis.} VIX/VIX3M backwardation (spot above three-month) marks stress that mean-reverts, so a high ratio predicts higher returns. \textbf{Data.} 1,234 aligned days, 2021-04-01 to 2026-03-02.\src{data/research/2026-04-01\_vix\_term\_structure.md}

\begin{table}[htbp]
\centering\small
\caption{VIX/VIX3M ratio against 21-day SPY forward returns.}
\begin{tabular}{lrrr}
\toprule
Sample & Spearman & $t$ & $p$ \\
\midrule
In-sample (863 rows) & +0.198 & 5.94 & $<0.001$ \\
Out-of-sample (371 rows) & +0.176 & 3.44 & 0.0006 \\
2023 only & $-0.019$ & & 0.763 \\
2026 to April & $-0.501$ & & 0.001 \\
\bottomrule
\end{tabular}
\end{table}

Quintile 5 (backwardation) returned 2.58\% per 21 days against 0.64\% for quintile 1, and backwardation occurred on 4.4\% of days. An independent reproduction matched every headline figure within 5\% and added the year-by-year split: the signal is dead in 2023 and reversed in 2026. The ratio's correlation with the VIX level is 0.74 (Pearson); after residualizing on VIX level the incremental Spearman is 0.12. On non-overlapping windows ($n = 59$) the Spearman is 0.33 with $t = 2.60$. \textbf{Verdict.} Conditional: regime-gated overlay only, zero weight in risk-on, capped at 5 to 8\%. The production configuration that followed gave it 0.05 in risk-on and zero elsewhere, the opposite of the conditions. It was removed from the composite as zero-contribution and is not deployed.

\subsection{Put/call ratio (proxies)}

\textbf{Hypothesis.} Extreme put buying is contrarian bullish. The CBOE series is not available through the data pipeline, so a 252-day VIX percentile and a 63-day SKEW z-score were tested as proxies over 943 days.\src{data/research/2026-04-01\_put\_call\_ratio.md} In-sample $p$-values ran from 0.15 to 0.96. Quintiles were not monotonic: the calm quintiles returned 1.46\% and 1.60\% per 21 days, the middle 0.36\% and 0.45\%, the fear quintile 2.32\%. Out-of-sample significance appeared where in-sample had none, which is the signature of a period effect. \textbf{Verdict.} Rejected.

\subsection{Insider transactions}

\textbf{Hypothesis.} Net insider buying predicts returns. \textbf{Data.} 801 Form 4 records across the nine names, 2024-04 to 2026-03, aggregated monthly to 23 observations.\src{data/research/2026-04-01\_insider\_trading.md} Purchases were 6 of 801; sales 413. No aggregate or per-ticker test reached $p < 0.05$; the average per-ticker Spearman was $-0.02$ and signs were split four to four. \textbf{Verdict.} Rejected: the universe is the wrong one for the signal and the data source has two years of history.

\subsection{Credit spreads}

\textbf{Hypothesis.} Widening high-yield versus investment-grade spreads predict lower equity returns. Proxied by HYG and LQD returns over 1,540 days from 2020-01-14.\src{data/research/2026-04-01\_credit\_spreads.md}

\begin{table}[htbp]
\centering\small
\caption{Credit-ratio z-score against 21-day portfolio forward returns.}
\begin{tabular}{lrrr}
\toprule
Sample & Spearman & $t$ & $p$ \\
\midrule
In-sample (1,078 rows) & $-0.195$ & $-5.48$ & $<0.001$ \\
Out-of-sample (462 rows) & +0.021 & +0.59 & 0.553 \\
\bottomrule
\end{tabular}
\end{table}

The 21-day spread variant flipped sign between samples ($+0.06$ to $-0.19$). Rolling ICIRs were $-0.53$ to $+0.34$. Turnover of the z-score implied 5\% a year in costs at 5 basis points. \textbf{Verdict.} Rejected.

\subsection{Sector relative strength}

\textbf{Hypothesis.} Names in strong sectors outperform (industry momentum). Eleven sector ETFs mapped to the nine names, 1,540 days.\src{data/research/2026-04-01\_sector\_relative\_strength.md} No in-sample result survived Bonferroni (best $p = 0.013$ against 0.0083). All twelve out-of-sample Pearson correlations were negative (eleven of twelve Spearman), and the cross-sectional variant was significantly negative ($p = 0.014$). Every rolling ICIR was negative. Quintiles were not monotonic. \textbf{Verdict.} Rejected.

\subsection{Breadth}

Two constructions were tested: the share of the nine names above their 50-day average, and 21-day momentum of the RSP/SPY ratio.\src{data/research/2026-04-01\_breadth\_indicators.md} The first had a raw rolling ICIR of $-1.45$ and passed Bonferroni out of sample on all four targets, but correlated 0.77 with the trend signal and 0.68 with RSI. A multivariate follow-up regressed forward returns on trend, RSI, the VIX term structure and each breadth variant.\src{data/research/2026-04-01\_breadth\_multivariate.md}

\begin{table}[htbp]
\centering\small
\caption{Marginal contribution of breadth variants after controlling for existing signals (21-day forward returns).}
\begin{tabular}{llrrrr}
\toprule
Variant & Sample & $\Delta R^2$ & $t$ & Marginal Spearman & VIF \\
\midrule
Share above 50d & IS & 0.0253 & $-4.87$ & $-0.157$ & 4.20 \\
Share above 50d & OOS & 0.0002 & +0.26 & $-0.027$ & \\
RSP/SPY 21d momentum & IS & 0.0044 & $-2.02$ & $-0.089$ & 1.14 \\
RSP/SPY 21d momentum & OOS & 0.1463 & $-8.42$ & $-0.428$ & \\
\bottomrule
\end{tabular}
\end{table}

The level variant is redundant. The momentum variant's explanatory power rose 33-fold from in-sample to out-of-sample, the trend coefficient flipped sign between samples, and the overlap-adjusted out-of-sample $t$ is about 1.84. \textbf{Verdict.} Both rejected.

\subsection{Options skew}

\textbf{Hypothesis.} A high CBOE SKEW index is contrarian bullish. Six constructions over 1,202 days.\src{data/research/2026-04-01\_options\_skew.md} No variant passed Bonferroni in-sample. The percentile variant went from $+0.075$ in-sample to $-0.252$ out-of-sample. Quintiles were hump-shaped. The z-score variant implied 6.3\% a year in turnover costs at 5 basis points. \textbf{Verdict.} Rejected.

\subsection{Short interest}

Only current snapshots exist in the pipeline, giving one cross-section of nine names, all with short interest between 0.85\% and 1.42\% of float.\src{data/research/2026-04-01\_short\_interest.md} A volume-spike proxy had all $p > 0.2$. \textbf{Verdict.} Rejected for lack of data.

\subsection{Variance risk premium}

\textbf{Hypothesis.} Implied minus realized volatility predicts positive returns (Bollerslev, Tauchen and Zhou, 2009). 1,375 days.\src{data/research/2026-04-01\_variance\_risk\_premium.md} The 21-day IC was $+0.067$ in-sample ($p = 0.047$) and $-0.078$ out-of-sample. In the low-VIX regime the IC was $-0.38$ ($p < 0.001$), the opposite of the hypothesis. The premium was positive on 87.4\% of days, so it rarely carries timing information. \textbf{Verdict.} Rejected.

\subsection{Bond-equity correlation regime}

\textbf{Hypothesis.} The z-score of the 63-day SPY/TLT correlation predicts returns. 1,375 days.\src{data/research/2026-04-01\_bond\_equity\_correlation.md} The 21-day IC was $+0.101$ in-sample ($p = 0.003$) and $-0.136$ out-of-sample ($p = 0.007$): significant in both samples with opposite signs. Per-stock results split growth (GOOGL $-0.39$, AMZN $-0.39$) from defensives (JNJ $+0.32$, KO $+0.21$), which suggests a rotation indicator rather than a timing one. \textbf{Verdict.} Rejected.

\subsection{Copper/gold ratio}
\label{sec:cuau}

\textbf{Hypothesis.} Rising copper relative to gold signals growth. The data said the opposite: the 21-day ratio momentum had a negative IC against VOO forward returns in both samples ($-0.099$ in, $-0.256$ out at 21 days), monotonic out-of-sample quintiles (2.98\% to 0.18\%), and survived Bonferroni over nine tests at 0.1\%.\src{data/research/2026-04-01\_copper\_gold\_ratio.md} The sign flips with regime: $+0.29$ when VIX is below 15, $-0.43$ above 25. A walk-forward with the inverted sign fixed in advance followed.\src{data/research/2026-04-01\_copper\_gold\_walk\_forward.md}

\begin{table}[htbp]
\centering\small
\caption{Copper/gold walk-forward, 72 monthly periods 2020-02 to 2026-02, 15 bps per side.}
\begin{tabular}{lr}
\toprule
Mean walk-forward IC & 0.155 \\
Walk-forward ICIR & 0.355 \\
Periods with positive IC & 66.7\% \\
$t$-test on period ICs & $t = 3.01$, $p = 0.0036$ \\
Bonferroni over 14 signal families & $p = 0.050$ \\
Bonferroni over about 42 sub-tests & $p = 0.150$ \\
Long-short Sharpe, net & 0.98 \\
Crisis ICIR (16 periods) & 1.29 \\
Risk-on ICIR (9 periods) & $-0.18$ \\
First-half vs second-half IC & 0.192 vs 0.119 \\
IC retained after dollar orthogonalization & 98.9\% \\
\bottomrule
\end{tabular}
\end{table}

\textbf{Verdict.} Conditional. It was implemented with regime gating, then removed from the composite when the null experiment (Section~\ref{sec:null}) showed the full stack indistinguishable from equal weight. Not deployed.

\subsection{Earnings revision momentum and analyst dispersion}

Price-target revisions over 30, 60 and 90 days had positive in-sample ICs (0.04 to 0.09) and negative out-of-sample ICs of the same size across all horizons; the 60-day revision was $+0.086$ then $-0.086$.\src{data/research/2026-04-01\_earnings\_revision\_momentum.md} An upgrade-minus-downgrade count had a consistent sign (0.055 to 0.100) and a rolling ICIR of 0.59, but was never walk-forward tested and correlates 0.41 to 0.69 with price momentum. Analyst dispersion could not be backtested at all: one cross-section of nine targets.\src{data/research/2026-04-01\_analyst\_estimate\_dispersion.md} \textbf{Verdict.} Both rejected.

\subsection{Earnings surprise (PEAD)}

The original engine carried a post-earnings-drift signal built on Yahoo earnings-surprise fields. The fields are sparse for mega-caps and the walk-forward assigned the signal no weight in any regime; it was removed.\src{quantshield/config.py} \textbf{Verdict.} Rejected.

\subsection{India: currency momentum, foreign flows}

USDINR momentum and z-score variants flipped sign between samples (21-day Spearman $+0.135$ in-sample, $-0.219$ out-of-sample for the z-score), the exporter-minus-importer spread went the wrong way out of sample ($-0.30$, $p < 0.001$), and after residualizing on UUP momentum the incremental IC was $-0.027$ ($p = 0.37$).\src{data/research/2026-04-01\_currency\_momentum\_inr.md} Foreign flow was proxied by INDA volume and the INDA/EEM price ratio. Direct flow momentum flipped sign ($+0.098$ to $-0.159$). The ratio's 20-day momentum had an out-of-sample IC of 0.242 against 21-day Nifty returns and survived Bonferroni over 27 tests, but its rolling ICIR was $-0.06$ and it reversed when India VIX exceeded 25 (IC $-0.30$).\src{data/research/2026-04-01\_fii\_dii\_flow\_momentum.md} A follow-up with 27 further formulations found four significant at 10\% against 2.7 expected by chance.\src{data/research/2026-04-01\_fii\_flow\_deep\_dive.md} \textbf{Verdict.} Rejected; the ratio was never walk-forward tested.

\subsection{Seasonality and turn-of-month}

Fifteen years of SPY and Nifty daily data, four calendar effects.\src{data/research/2026-04-01\_seasonality\_signals.md} In the US every effect had $p > 0.35$ and the composite IC was 0.02 at five days. In India the turn-of-month window (last three and first three trading days) returned 25.18\% annualized against 4.36\% outside it, a spread of 8.26 basis points per day, $p = 0.011$ full-sample and $p = 0.013$ out-of-sample, against a Bonferroni threshold of 0.0033. \textbf{Verdict.} US rejected. India turn-of-month is conditional: an operational rule about when to rebalance, not a composite weight, and not enforced in code.

\subsection{Regime detection with a hidden Markov model}
\label{sec:hmm}

A Gaussian HMM on VIX level and change, 21-day benchmark return and volatility, yield change and gold return, 2,471 days from 2016.\src{data/research/2026-04-01\_hmm\_regime\_detection.md} The three-state fit collapsed two states onto each other (mean VIX 14.6 and 14.7) and alternated between them daily (transition probabilities 0.999 and 0.970), producing 156.5 transitions a year against 22.8 for the threshold heuristic. The heuristic's crisis label (141 days) preceded a mean 21-day return of $+5.74\%$; the HMM's (954 days) $+1.63\%$. Walk-forward cumulative return was 179.0\% for the HMM and 177.1\% for the heuristic. \textbf{Verdict.} Rejected as a production component. A companion study in India used the heuristic to time equity, gold and cash and found crisis flags followed by a mean $+1.16\%$ five-day return, that is, the flag is contrarian.\src{data/research/multi-asset-regime-timing.md}

\subsection{MOVE/VIX divergence}

The 63-day z-score of the MOVE/VIX ratio over 2,827 days had an IC of $-0.017$ ($p = 0.37$) at five days and $+0.007$ ($p = 0.71$) at 21 days.\src{data/research/2026-04-02\_move\_vix\_divergence.md} \textbf{Verdict.} Rejected.

\subsection{Factor crowding and the momentum lookback}

A long-short 12-1 momentum portfolio in the nine names had a Sharpe of 0.164 over 2021 to 2026, worst months of $-15.2\%$ (January 2023) and $-12.9\%$ (February 2025), and the equal-weight universe itself correlated 0.77 with MTUM and 0.79 with SPMO.\src{data/research/2026-04-02\_factor\_crowding.md} The same study reported a 3-1 lookback Sharpe of 0.864 on the five-year window. On ten years the figures are 0.158 (3-1) and 0.225 (12-1); the sweep in Section~\ref{sec:lookback} closes the question.

\newpage
\section{Validation Results}

\subsection{Engine walk-forward}
\label{sec:wf}

Table~\ref{tab:wf} reports the \wfRunDate{} run\src{dashboard/src/data/us.json, dashboard/src/data/india.json} of the pipeline in Section~1: the five-signal composite, the HRP blend, limits and CVaR, no vol targeting, a flat 10 basis points for the US and the CNC schedule for India. The expanding window uses a \wfMinTrain-day minimum training set and \wfStep-day steps. No step fell back to equal weight in either market (\usWFfallbacks{} and \inWFfallbacks{} fallback periods).

\begin{table}[htbp]
\centering\footnotesize
\caption{Walk-forward summary, run of \wfRunDate. Return, alpha and drawdown are cumulative percentages over the window.}
\label{tab:wf}
\begin{tabular}{lrr}
\toprule
 & US vs VOO & India vs Nifty 50 \\
\midrule
Window & \usWFstart{} to \usWFend & \inWFstart{} to \inWFend \\
Periods (21 trading days) & \usWFperiods & \inWFperiods \\
Cost per rebalance & 10 bps one-way turnover & CNC schedule per leg \\
Portfolio return & \usWFport\% & \inWFport\% \\
Benchmark return & \usWFbench\% & \inWFbench\% \\
Excess return & \usWFalpha{} pp & \inWFalpha{} pp \\
Sharpe, portfolio / benchmark & \usWFsharpe{} / \usBenchSharpe & \inWFsharpe{} / \inBenchSharpe \\
Volatility, portfolio / benchmark & \usWFvol\% / \usBenchVol\% & \inWFvol\% / \inBenchVol\% \\
Max drawdown, portfolio / benchmark & \usWFmaxdd\% / \usBenchMaxdd\% & \inWFmaxdd\% / \inBenchMaxdd\% \\
Periods with positive alpha & \usWFwins{} (\usWFwinrate\%) & \inWFwins{} (\inWFwinrate\%) \\
Period-alpha $t$ / $p$ & \usWFtstat{} / \usWFpvalue & \inWFtstat{} / \inWFpvalue \\
Bootstrap 95\% Sharpe interval & $[\usSharpeCIlo, \usSharpeCIhi]$ & $[\inSharpeCIlo, \inSharpeCIhi]$ \\
Bootstrap 95\% alpha interval, pp/yr & $[\usAlphaCIlo, \usAlphaCIhi]$ & $[\inAlphaCIlo, \inAlphaCIhi]$ \\
Deflated Sharpe $p$ with $N = \ledgerRows$ & \usDSRpvalue & \inDSRpvalue \\
In-sample 252-day Sharpe, portfolio / benchmark & \usISsharpe{} / \usISbenchSharpe & \inISsharpe{} / \inISbenchSharpe \\
\bottomrule
\end{tabular}
\end{table}

Neither alpha interval excludes zero. The India $t$ of \inWFtstat{} ($p = \inWFpvalue$) is the closest any benchmark-relative test comes to the 5\% line; it is one of \ledgerRows{} trials and its alpha interval still includes zero. The US in-sample Sharpe of \usISsharpe{} over the trailing year against a walk-forward Sharpe of \usWFsharpe{} over six years is the usual gap between a fitted window and an honest one. The India in-sample return over the trailing year was \inISport\% against \inISbench\% for the benchmark.

By regime, the US walk-forward alpha came from risk-on months: \usRegimeOnMonths{} months with \usRegimeOnAlpha{} points of cumulative alpha and a \usRegimeOnWin\% win rate, against \usRegimeOffMonths{} risk-off months with \usRegimeOffAlpha{} points and \usRegimeOffWin\%, and \usRegimeCrisisMonths{} crisis months with \usRegimeCrisisAlpha{} points and \usRegimeCrisisWin\%. India was positive in every regime but thinly: \inRegimeOnMonths{} risk-on months with \inRegimeOnAlpha{} points (\inRegimeOnWin\%), \inRegimeOffMonths{} risk-off months with \inRegimeOffAlpha{} points (\inRegimeOffWin\%), and \inRegimeCrisisMonths{} crisis months with \inRegimeCrisisAlpha{} points.\src{regime\_performance in dashboard/src/data/us.json and india.json} Six and two crisis months are not samples, and the label is the regime the detector held at the rebalance, not the one the market turned out to be in.

\subsection{Deflated Sharpe ratio}
\label{sec:dsr}

Table~\ref{tab:dsr} applies the deflated Sharpe ratio of Section~1 to the daily walk-forward series of both markets with $N = \ledgerRows$.\src{deflated\_sharpe in dashboard/src/data/us.json and india.json}

\begin{table}[htbp]
\centering\small
\caption{Deflated Sharpe ratio on the walk-forward daily series, run of \wfRunDate.}
\label{tab:dsr}
\begin{tabular}{lrr}
\toprule
 & US & India \\
\midrule
Observed Sharpe, annualized & \usDSRobserved & \inDSRobserved \\
Benchmark Sharpe, annualized & \usDSRbenchmark & \inDSRbenchmark \\
Expected maximum null Sharpe, $N = \ledgerRows$ & \usDSRemax & \inDSRemax \\
Daily observations & \usDSRtobs & \inDSRtobs \\
Skewness / excess kurtosis & \usDSRskew{} / \usDSRkurt & \inDSRskew{} / \inDSRkurt \\
Probabilistic Sharpe ratio & \usDSRpsr & \inDSRpsr \\
$p$ & \usDSRpvalue & \inDSRpvalue \\
\bottomrule
\end{tabular}
\end{table}

The expected maximum Sharpe that \ledgerRows{} null trials would produce by chance is \usDSRemax{} in the US, just under the observed \usDSRobserved, and \inDSRemax{} in India, above the observed \inDSRobserved. $\mathrm{SR}^{*}$ binds on the expected maximum in both markets, since both benchmark Sharpes sit below it. The probabilistic Sharpe ratios are \usDSRpsr{} and \inDSRpsr, so $p = \usDSRpvalue$ and $p = \inDSRpvalue$. Neither survives deflation. Both $p$-values are lower bounds twice over: $N$ counts ledger rows rather than every variant tried, and the variance in the expected-maximum formula is the estimation variance of one Sharpe rather than the spread across trials.

\subsection{Factor decomposition}
\label{sec:ff}

The US walk-forward daily series was compounded into \ffMonths{} calendar months and regressed on the six Ken French factors with \ffLags{} Newey-West lags.\src{fama\_french in dashboard/src/data/us.json}

\begin{table}[htbp]
\centering\small
\caption{Six-factor regression of monthly US walk-forward excess returns, \ffMonths{} months, Newey-West with \ffLags{} lags.}
\label{tab:ff}
\begin{tabular}{lrrr}
\toprule
 & Coefficient & $t$ & $p$ \\
\midrule
Intercept, \% per month & \ffAlphaMonthly & \ffAlphaT & \ffAlphaP \\
Mkt-RF & \ffMktBeta & \ffMktT & \ffMktP \\
SMB & \ffSmbBeta & \ffSmbT & \ffSmbP \\
HML & \ffHmlBeta & \ffHmlT & \ffHmlP \\
RMW & \ffRmwBeta & \ffRmwT & \ffRmwP \\
CMA & \ffCmaBeta & \ffCmaT & \ffCmaP \\
Mom & \ffMomBeta & \ffMomT & \ffMomP \\
\midrule
$R^2$ & \ffRsq & & \\
\bottomrule
\end{tabular}
\end{table}

The intercept is \ffAlphaMonthly\% a month, \ffAlphaAnnual{} points a year, with $t = \ffAlphaT$ and $p = \ffAlphaP$. The fit ($R^2 = \ffRsq$) is carried by a market beta of \ffMktBeta{} and a profitability loading (RMW) of \ffRmwBeta{} with $t = \ffRmwT$; the momentum loading is \ffMomBeta{} and not distinguishable from zero, which matches the null experiment's finding that the momentum tilt does not move the portfolio. Three things keep the intercept from being a claim of alpha. It is one test among \ledgerRows{} and is not deflated. The universe was chosen in 2026 from names that are profitable large caps today, which is what a positive RMW loading and a positive intercept on \ffMonths{} months of their returns describe; the survivorship estimate of Section~\ref{sec:survivorship}, $+4.1$ points a year, is of the same order as the intercept. And the same series measured against its own benchmark, VOO, carries $p = \usWFpvalue$. No factor regression exists for India: the pipeline has no Indian factor library, and none is claimed.

\subsection{Bootstrap and parameter robustness}

A separate robustness run varied the momentum lookback (126 to 378 days) and the rebalance step (10, 21, 42 days) and perturbed every regime weight by $\pm 20\%$ in 100 Monte Carlo trials.\src{data/research/wf-robustness-analysis.md} Cumulative US alpha across the nine-cell grid ranged from 54\% to 92\%, always positive, never significant ($p$ from 0.07 to 0.20). India ranged from 60\% to 86\% and was significant in every cell ($p$ from 0.007 to 0.038) in that run. Weight perturbation moved US alpha by a standard deviation of 1.31 points around 81.6 with no sign flips. The direction is stable; the magnitude is not distinguishable from noise in the US.

\subsection{Equal-weight null experiment}
\label{sec:null}

Three strategies over 107 monthly periods from 2016-04 to 2026-04 in the US universe, all built on one code path with identical daily tracking, identical Sharpe arithmetic and identical cost function.\src{data/research/existential-experiment-v2.md} A: naive equal weight, monthly reset. B: equal weight followed by vol targeting, sector, single-name, beta and CVaR limits, and regime-adaptive rebalancing. C: the full engine.

\begin{table}[htbp]
\centering\small
\caption{Null experiment v2, US, 5 bps spread plus square-root impact.}
\begin{tabular}{lrrr}
\toprule
 & A: equal weight & B: EW + overlay & C: full engine \\
\midrule
Sharpe & 1.2716 & 1.2497 & 1.2473 \\
Annual return & 25.71\% & 22.43\% & 22.47\% \\
Annual volatility & 20.22\% & 17.95\% & 18.01\% \\
Max drawdown & $-27.35\%$ & $-26.10\%$ & $-25.89\%$ \\
Monthly turnover & 0.00\% & 1.64\% & 8.26\% \\
\bottomrule
\end{tabular}
\end{table}

Bootstrap 95\% intervals on the Sharpe difference: B minus A $[-1.22, +1.10]$, C minus A $[-0.13, +0.19]$, C minus B $[-1.07, +1.24]$. Paired $t$ on monthly returns: A versus B $p = 0.34$, A versus C $p = 0.045$. The ranking held at 10 basis points flat and at zero cost. The paired $t$ on A versus C tests mean return, which is lower for C by construction because the overlay reduces exposure; the Sharpe interval is the relevant test and it includes zero. An earlier version of this experiment reported C at 0.98 because C's Sharpe came from a different calculation path and its daily series was synthesized from monthly returns; that figure is withdrawn. The India run of the earlier version, whose A and B legs share one code path, gave A 1.20 and B 1.13 with a bootstrap interval of $[-0.12, +0.05]$ over 108 periods.\src{data/research/existential-experiment.md} A review of the earlier version also withdrew a recovery-time claim: on the same event (COVID) equal weight recovered in 67 trading days and the overlay in 84.

\subsection{Survivorship bias}
\label{sec:survivorship}

Equal-weight portfolios were built from the current universe, from a plausible 2021 universe (AAPL, MSFT, AMZN, GOOGL, BRK-B, JNJ, XOM, WFC, T), and from the current universe plus six fallen large caps (GE, XOM, WFC, IBM, T, INTC), all over 2021-04 to 2026-04.\src{data/research/2026-04-02\_survivorship\_bias.md}

\begin{table}[htbp]
\centering\small
\caption{Survivorship estimate, five years, equal weight.}
\begin{tabular}{lrrrr}
\toprule
Portfolio & Sharpe & Annual return & Total & Max DD \\
\midrule
Current 9 (US) & 1.154 & 21.4\% & 166.3\% & $-27.4\%$ \\
Historical 9 (2021 view) & 1.078 & 17.3\% & 121.9\% & $-18.9\%$ \\
Current plus 6 fallen & 1.263 & 21.2\% & 167.1\% & $-23.9\%$ \\
VOO & 0.741 & 12.5\% & 73.2\% & $-24.5\%$ \\
Current 20 (India) & 1.056 & 14.6\% & 95.3\% & $-15.2\%$ \\
Historical 20 (India) & 0.898 & 13.1\% & 80.3\% & $-17.1\%$ \\
\bottomrule
\end{tabular}
\end{table}

The gap is $+4.1$ percentage points a year in the US and $+1.5$ in India. Two former Nifty constituents (DHFL, Tata Motors in its earlier form) could not be downloaded at all, so the India figure is a floor. The walk-forward removes look-ahead from signal values but not from the choice of names; every alpha in this report sits inside these margins.

\subsection{Momentum lookback sweep}
\label{sec:lookback}

Lookbacks of 2, 3, 4, 6, 9 and 12 months, each skipping the latest month, tilting an equal-weight base at 50\% in walk-forward with a 252-day minimum window, 2015 to 2026.\src{data/research/momentum-lookback-study.md}

\begin{table}[htbp]
\centering\footnotesize\setlength{\tabcolsep}{4pt}
\caption{Lookback sweep. Alpha is net annual alpha; P is the bootstrap probability that the variant beats 12-1 on Sharpe.}
\begin{tabular}{lrrrrrr}
\toprule
Lookback & US Sharpe & US alpha & P & India Sharpe & India alpha & India ICIR \\
\midrule
2-1 & 1.342 & 8.37\% & 46.8\% & 1.180 & 7.78\% & $-0.03$ \\
3-1 & 1.332 & 8.32\% & 46.1\% & 1.157 & 7.45\% & $-0.07$ \\
4-1 & 1.329 & 8.40\% & 47.4\% & 1.158 & 7.54\% & $-0.06$ \\
6-1 & 1.378 & 9.65\% & 50.5\% & 1.188 & 8.20\% & 0.07 \\
9-1 & 1.364 & 9.52\% & 50.0\% & 1.191 & 8.21\% & 0.01 \\
12-1 & 1.366 & 9.39\% & n/a & 1.232 & 8.80\% & 0.05 \\
\bottomrule
\end{tabular}
\end{table}

No variant is distinguishable from 12-1. All ICIRs are below 0.15, so the tilt works through portfolio-level exposure rather than stock selection. An independent ten-year run embedding 24 lookback and skip variants in the HRP-based pipeline produced walk-forward Sharpes between 0.937 and 1.036, with equal weight and no tilt at 1.011 and the production 12-1 tilt at 0.966. The lookback is a second-order parameter and the tilt does not beat equal weight.

\subsection{Sector equal weight}

Equal weight across sectors, then equal within each, was tested as a construction default.\src{data/research/sector-equal-weight-study.md; data/research/sector-ew-deep-validation.md} Full-sample India Sharpe was 1.218 against 1.056 for naive equal weight in the first study (1.223 against 1.074 in the deep validation). The walk-forward difference was $+0.066$ with a bootstrap interval of $[-0.121, +0.218]$ over 75 periods. Random sector assignments with the same bucket counts produced differences between $-0.02$ and $+0.02$, so the effect is the specific India map (single-stock sectors that happened to outperform banks), not the mechanism. In the US the same construction cut Sharpe by 0.104. Only 59.4\% of rolling two-year windows favored it. Costs do not separate the two (turnover 5.02\% vs 4.95\% a month). \textbf{Verdict.} A zero-cost default with a concentration rationale, not a source of excess return. The engine uses sector caps instead.

\subsection{Portfolio construction comparison}

Equal weight, inverse volatility, equal risk contribution and HRP, all without signal tilts, in walk-forward with 10 basis point costs: Sharpe 1.293, 1.236, 1.272 and 1.243.\src{data/research/2026-04-02\_portfolio\_construction\_comparison.md} The spread of 0.056 is inside estimation error at nine names. HRP was kept because it was already implemented and does no harm.

\subsection{Alternative architectures}

Four re-configurations were tested against equal weight on the US universe over 2016 to 2026.\src{data/research/signal-rethink.md}

\begin{table}[htbp]
\centering\small
\caption{Alternative architectures against equal weight (Sharpe 1.272, drawdown $-27.35\%$, recovery 296 days).}
\footnotesize
\begin{tabular}{lrrrr}
\toprule
Strategy & Sharpe & Difference & Bootstrap CI & Verdict \\
\midrule
Regime-based cash allocation & 1.220 & $-0.052$ & $[-0.30, +0.15]$ & rejected \\
Sector ETF universe, equal weight & 0.654 & $-0.618$ & & rejected \\
Sector ETF universe, momentum tilt & 0.703 & $-0.569$ & & rejected \\
Crash buying with 20\% leverage & 1.179 & $-0.093$ & $[-0.12, +0.02]$ & rejected \\
Crash buying from 20\% cash buffer & 1.206 & $-0.066$ & $[-0.09, +0.06]$ & rejected \\
\bottomrule
\end{tabular}
\end{table}

The cash-buffer variant cut the worst recovery from 296 days to 73 at a cost of 4.5 points of annual return; that is a preference trade, not alpha. Country ETF rotation (six ETFs, Sharpe 0.993 full-sample, 0.583 walk-forward, interval against equal weight $[-0.67, +1.49]$), NSE midcap momentum (equal weight 1.214 against the best momentum variant 0.964, every interval including zero) and a contribution-timing rule (no US timing value; India $+10.1\%$ over 18 years on two episodes) were also rejected.\src{data/research/country-rotation-study.md, midcap-momentum-study.md, dca-crash-buy-timing.md}

\subsection{Stress replays}

Walk-forward replays of the composite through three episodes, 21-day rebalances, 10 basis point costs.\src{data/research/stress-test-results.md}

\begin{table}[htbp]
\centering\small
\caption{Stress replays.}
\footnotesize
\begin{tabular}{lrrrr}
\toprule
Episode & Engine & Benchmark & Engine max DD & Benchmark max DD \\
\midrule
COVID, US, Jan to Jun 2020 & +12.78\% & $-4.08\%$ & $-25.89\%$ & $-33.99\%$ \\
Rate hikes, US, 2022 & $-13.36\%$ & $-18.67\%$ & $-22.35\%$ & $-24.52\%$ \\
COVID, India, Jan to Jun 2020 & $-11.78\%$ & $-16.12\%$ & $-35.80\%$ & $-38.44\%$ \\
\bottomrule
\end{tabular}
\end{table}

The detector first labeled crisis on 2020-03-04 at VIX 32, nine trading days after the peak, and held it through May. The 2020 outperformance came from the crisis-regime mean-reversion weight buying oversold names into a V-shaped recovery. In 2022 the regime changed six times in twelve months and the outperformance came from HRP's tilt toward low-volatility names, not from signals. The India replay used the US detector and is a weaker test. These are three episodes from one decade; they show behavior, not expected value.

\newpage
\section{Live Execution and Safety Model}

The only book that trades real money is a small India account on Zerodha, run by three modules.\src{quantshield/live/planner.py, quantshield/live/executor.py, quantshield/live/daemon.py}

\textbf{Planner.} Holds NIFTYBEES as a 50\% core and two satellite slots filled from the India composite (same five signals, same regime weights). A held satellite is kept while it ranks in the top four; one that falls outside the top four is replaced only by a candidate whose score exceeds its own by at least 0.15; one that sits below rank 12 for three sessions is sold; a 20\% loss against average cost forces a sale; a 50\% look-through cap on financials (counting 35\% of the core as financial) blocks buys; trades below 800 INR are not planned; a suspected corporate action suppresses exits for five sessions. The planner writes a plan file and never places orders.

\textbf{Executor.} Reads a plan generated the same day and no older than six hours, reconciles broker holdings against local state and aborts on mismatch, then places LIMIT CNC DAY orders at a 1\% band, modifies once to a 2\% band after 90 seconds, and cancels after a further 60 seconds. Caps: 3,000 INR per order, 6,000 INR turnover per day, six orders per day. Every order is journaled before placement so a retry cannot duplicate it. If the placement call fails ambiguously and the follow-up scan of the order book also fails, the leg is journaled as \code{PLACEMENT\_UNVERIFIED} and a notification is sent; nothing is retried. A \code{data/KILL} file halts execution before and between orders. A file lock prevents two instances. Live mode requires both \code{AUTO\_EXECUTE} and \code{ZERODHA\_LIVE\_MODE} to be set and \code{-{}-dry-run} absent; the default is a dry run. Fills, limit prices and slippage in basis points are appended to a trade log.

\textbf{Daemon.} During the trading window it generates the plan if none exists for the day, runs the executor once, takes a daily NAV snapshot after the close, and checks feed freshness. Every step requires a Kite access token dated today, which comes from a manual login; when the token is stale the daemon sends one notification and does nothing, and the snapshot step applies the same check. While \code{data/KILL} exists the daemon does not start the executor and says so once a day. An intraday feed heartbeat older than a day is ignored rather than reported as stale, and a failed dashboard deploy is reported once a day. The NSE holiday table it consults carries all sixteen weekday closures of 2026.\src{quantshield/calendar.py}

\textbf{Costs.} The planner and executor price each leg with the delivery cost model of Section~1 (STT, exchange and SEBI charges, GST, stamp duty, depository charge, ETF exemptions) and scale buy quantities down until the order plus its cost fits available cash.

\textbf{Record.} A daily snapshot file records portfolio value, cash and the value of the NIFTYBEES units the starting capital would have bought, so the benchmark is a same-capital buy-and-hold of the core ETF from inception (2026-07-20). That file is runtime state and is not in the repository. The public export indexes both series to 100, publishes weights rather than amounts, and carries no cash figures, cost basis, client identifiers or token state.\src{quantshield/live/export.py, dashboard/public/live/dashboard.json}

\textbf{Status.} The public export of \wfRunDate{} holds 13 daily snapshots from 2026-07-20 to 2026-08-05: portfolio $+2.89\%$ against $+1.27\%$ for the benchmark, maximum drawdown $-2.40\%$, three positions (NIFTYBEES 49.6\%, BAJFINANCE 21.1\%, SBIN 20.9\% of account value), zero fills, and a last plan dated 2026-07-21 with no orders and one warning, the financials look-through above the 50\% cap.\src{dashboard/public/live/dashboard.json} The opening positions were placed by hand, outside the executor. No execution journal entry exists. The loop has never placed an order: on the trading days since inception the daily token precondition was not met, so the executor did not run, and after 2026-08-05 the snapshot step stopped for the same reason. The guardrails are tested in code; they are untested against a live fill.

\newpage
\section{Pre-registered Intraday Study}
\label{sec:orb}

An opening-range breakout on NIFTYBEES was written down before any data was collected.\src{data/intraday/gate.json} Registered 2026-07-21: opening range from the 1-minute bars of 09:15 to 09:45 IST; one buy stop per session at range high plus one tick, valid until 13:00; initial stop at the range midpoint; exit on stop or at 15:10; 2,000 INR notional, long only, no re-entries; entries pay one tick plus half a spread, exits pay half a spread; intraday (MIS) costs from the cost module. The gate to any live pilot requires at least 60 live paper sessions and a one-sided bootstrap $p$ below 0.05 on the exposure-adjusted daily delta against same-capital buy-and-hold, cannot open in a window where raw buy-and-hold was negative, and freezes the rules: any parameter change starts a new arm with its own clock.

A replay on official 1-minute candles was run the same day as a plausibility check, explicitly outside the gate.\src{data/intraday/orb\_replay.json} The test statistic is the share of 10,000 resampled means of the exposure-adjusted daily delta that fall at or below zero, $P(\mathrm{mean} \le 0)$; the gate reads it as a one-sided $p$, and small values favour a positive mean.

\begin{table}[htbp]
\centering\small
\caption{ORB-NIFTYBEES-v1 replay, 2026-01-05 to 2026-07-20, after both amendments of \wfRunDate.}
\begin{tabular}{lr}
\toprule
Sessions & 132 \\
Triggered & 43 (32.6\%) \\
Wins & 13 (30.2\%) \\
Average win / average loss & 3.76 / $-7.03$ INR \\
Payoff ratio & 0.54 \\
Breakeven win rate & 65.1\% \\
Gross & $-76.75$ INR \\
Costs & 85.18 INR \\
Net & $-161.93$ INR on 2,000 notional ($-8.1\%$) \\
Max drawdown & $-167.37$ INR \\
Median opening-range width & 0.783\% (IQR 0.576\% to 1.124\%) \\
Median cost per unit of risk & 0.31 R \\
Benchmark, summed sessions / buy-and-hold & $+6.93\%$ / $-7.34\%$ \\
Exposure-adjusted delta, mean / sd & $-0.129\%$ / 0.202\% \\
Bootstrap $P(\mathrm{mean} \le 0)$, 10,000 resamples & 1.0 \\
\bottomrule
\end{tabular}
\end{table}

The first replay used the 09:15 candle open as the session benchmark open. That print is the pre-open auction and sits 0.35\% above the 09:16 open on average, so the summed benchmark return came out at $-38.63\%$ over a window whose first-close to last-close return was $-6.99\%$. A dated amendment of \wfRunDate{} in the gate file moved the benchmark open to the 09:15 bar close: the summed session benchmark became $+6.93\%$, the exposure-adjusted delta mean moved from $-0.0866\%$ to $-0.1214\%$, and $P(\mathrm{mean} \le 0)$ stayed at 1.0. The strategy-side figures never used that column and reproduced exactly. A second amendment of the same date records the cost correction of Section~1: the intraday model had charged no STT on ETF sells, and at the correct 0.025\% the costs on the same 43 round trips rose from 65.16 to 85.18 INR, net from $-141.91$ to $-161.93$ INR, the maximum drawdown from $-147.83$ to $-167.37$ INR, the breakeven win rate from 60.8\% to 65.1\%, and the exposure-adjusted delta mean to $-0.129\%$. Entries, exits and the gross loss of $-76.75$ INR did not change, and $P(\mathrm{mean} \le 0)$ is still 1.0. Same-bar stop ambiguity was resolved against the strategy. The live paper arm scored one session (2026-07-21, untriggered) against a 60-session requirement, because the feed job was never scheduled. The gate is closed and the strategy cannot beat its own cost geometry at this notional.

\newpage
\section{Known Limitations}

\begin{enumerate}
  \item \textbf{Survivorship.} Both universes were chosen with hindsight. The estimated inflation ($+4.1$ and $+1.5$ points per year) exceeds every alpha in Section~3. Point-in-time constituents were not available.
  \item \textbf{One regime sample.} The walk-forward window (2020 to 2026) contains one V-shaped crash, one grinding bear and a long recovery. The US crisis sub-sample is six months and the India one is two.
  \item \textbf{Nine names.} Cross-sectional signals have almost no dispersion to rank at $N = 9$, and equal weight has zero estimation error. Most academic anomalies are documented on hundreds of names.
  \item \textbf{Fixed hyperparameters.} Tilt strength, weight bounds, regime thresholds and regime weight tables were set before the walk-forward and are not re-estimated inside it.
  \item \textbf{Factor decomposition.} One regression on \ffMonths{} months of a hindsight universe. The intercept's $p$ of \ffAlphaP{} is not deflated for the trial count, and India has no factor data at all.
  \item \textbf{Deflated Sharpe.} $N = \ledgerRows$ counts ledger rows, not the parameter variants tried inside each study, and the variance input is the estimation variance of the observed Sharpe rather than the cross-trial variance, so the reported $p$ is a lower bound on two counts. The result is already negative at that floor ($p = \usDSRpvalue$ and \inDSRpvalue).
  \item \textbf{No live orders.} The India loop has a full guardrail set and a cost model checked against one contract note, but it has never fired because the daily token login is manual. The record is 13 snapshots and stops where the token did.
  \item \textbf{Conventions.} The within-window drift of the walk-forward, the risk-free rate (end-of-sample \code{\^{}TNX} for the US, a fixed 6.5\% for India), the Ward linkage in HRP, the one-day lag on the India macro inputs and the two beta windows (126 days for the cross-asset signal, 252 days for the beta cap) are choices, not estimates. Section~1 states each one; none was varied.
  \item \textbf{Intraday.} One live paper session against a 60-session gate. The replay is in-sample and single-regime; its benchmark column and its ETF sell tax were both wrong in the first replay and were corrected by dated amendments that left the verdict unchanged.
  \item \textbf{Overlapping returns.} Most signal studies use 21-day forward returns on daily rows; their $t$-statistics overstate significance by roughly $\sqrt{21}$ unless the study says otherwise.
  \item \textbf{Data source.} Yahoo Finance daily closes with no vendor guarantee; several studies were limited to proxies because the primary series was unavailable.
\end{enumerate}

\section{Conclusion}

Thirty-seven candidates were evaluated with a consistent protocol: split samples, information coefficients, quintiles, regime splits, cost estimates, correlation with existing signals, Bonferroni thresholds, and where a candidate survived, an expanding-window walk-forward with bootstrap intervals. Twenty-eight were rejected, most for sign flips between samples, redundancy with signals already in the composite, or data that did not exist in usable form. Three passed under conditions and are not deployed. Five remain in the composite, and a controlled comparison shows that the composite they form, with HRP and every risk limit applied, is statistically indistinguishable from putting equal money in the same nine names. The walk-forward excess returns of \usWFalpha{} and \inWFalpha{} cumulative points carry $p$-values of \usWFpvalue{} and \inWFpvalue{}, sit inside bootstrap intervals that include zero, fail the deflated Sharpe test at $N = \ledgerRows$, and annualize to about the survivorship inflation of the universes they were measured on.

That is the negative result. What the project does show is a research process that reached it: the tests were pre-specified often enough to catch a five-year momentum artifact, a calculation asymmetry that made signals look harmful, a recovery-time comparison across different events, and a sector-weighting effect that vanished under random relabeling. The intraday study was registered before its data existed and closed on its own terms. The live loop encodes every rule needed to trade a small book safely and a cost model checked against a broker note, and has not yet traded. The \wfRunDate{} run added the three results the previous draft still owed: a walk-forward on the current pipeline with CNC costs for India, a factor regression on real Ken French data, and a deflated Sharpe on the walk-forward series with $N = \ledgerRows$. All three point the same way as the rest. The one result still owed is a live fill.

\end{document}
